% wird automatisch alphabetisch sortiert
\newglossaryentry{frontend}{
  name={Frontend},
  description={Grafische Benutzeroberfläche zur Bedienung einer Website~\cite{an22}}
}
\newglossaryentry{backend}{
  name={Backend},
  description={Teil einer Website auf welche der Nutzer keinen Zugriff hat und dafür Dient Daten zu erfassen, verarbeiten und speichern~\cite{an23}}
}
\newglossaryentry{vektordatenbank}{
  name={Vektordatenbank},
  plural={Vektordatenbanken},
  description={In dieser Datenbank werden Datenpunkte als Zahlenreihen gespeichert, welche als Vektoren bezeihcnet werden und anhand der Ähnlichkeit verglichen und gruppiert werden~\cite{an24}}
}
\newglossaryentry{objectstorage}{
  name={Object Storage},
  description={Eine Datenarchitektur, die Informationen als Einheiten oder Objekte organisiert~\cite{an25}}
}
\newglossaryentry{nutzerdatenbank}{
  name={Nutzerdatenbanken},
  plural={Nutzerdatenbanken},
  description={Eine Datenbank, welche die Nutzerinformationen speichert~\cite{an26}}
}
\newglossaryentry{embeddingpipeline}{
  name={Embedding-Pipeline},
  description={Sie ermöglichen die transformation von Daten in Vektorräume, wodurch Beziehungen und Muster der Daten effizienter erkannt werden können~\cite{an27}}
}
\newglossaryentry{cicd}{
  name={CI/CD},
  description={Ein zweiteiliger Prozess, der die regelmäßige Integration, das Testen und die Bereitstellung der Codeänderungen umfasst~\cite{an28}}
}
\newglossaryentry{qdrant}{
  name={Qdrant},
  description={Eine Vektordatenbank, die Embeddings speichert und eine Simliarity Srach ermöglicht~\cite{an29}}
}